\clearpage\fancyhead[L]{\small JUNO.}\hypertarget{native-380-lower}{}\label{native-380-lower}
\begin{center}\textbf{JUNO.}\par\vspace{4pt}\hrule\vspace{4pt}\fontsize{10}{12}\selectfont Göttingische gelehrte Anzeigen. Stück 19. Seite 189. 1814 Jan. 31.\par\vspace{4pt}\hrule\vspace{4pt}\end{center}

\par\indent Noch fügen wir als Nachtrag zu der neulich gegebenen Nachricht über die Juno (m. s. St. 204 v. v. J.) drei Beobachtungen dieses Planeten von Hrn. Prof. \textsc{Bessel} bei:

\begin{center}\fontsize{10}{12}\selectfont\setlength{\tabcolsep}{3pt}\renewcommand{\arraystretch}{1.15}\begin{tabular}{l|c|rrr|c}
1812 & \shortstack{Mittlere Zeit\\in Königsberg} & \multicolumn{3}{c|}{\shortstack{Gerade\\Aufsteigung}} & \shortstack{Südliche\\Abweichung}\\\hline
Nov. 18 & $12^{\mathrm u}\;12^{\mathrm m}\;26^{\mathrm s}5$ & $60^\circ$ & $41'$ & $21''5$ & ---\quad---\\
21 & $11\;58\;20.4$ & 60 & 6 & 42.5 & ---\quad---\\
Dec. 11 & $10\;25\;54.9$ & 56 & 39 & 0.8 & $3^\circ\;53'\;27''2$\\
\end{tabular}\end{center}

\par\indent Dies sind die Erstlinge von der neuen Königsberger Sternwarte, deren Bau vor drei Jahren angefangen wurde, und welche Hr. \textsc{Bessel} schon im vorigen Herbst bezogen hat. Die Astronomie ist der Regierung den lebhaftesten Dank schuldig, die dieses für die Wissenschaft so wichtige Institut gegründet, mit vortrefflichen Instrumenten ausgerüstet, und selbst in diesen drangvollen Jahren so rasch vollendet hat: in der That wäre jeder Monat, um welchen ein Astronom wie \textsc{Bessel} später in seine volle Wirksamkeit gesetzt wäre, ein Verlust für die Wissenschaft gewesen.
